\documentclass[11pt,a4paper]{article}
\usepackage[margin=23mm]{geometry}
\usepackage[T1]{fontenc}
\usepackage{lmodern,xurl,hyperref}
\hypersetup{hidelinks}
\DeclareUnicodeCharacter{00B2}{\ensuremath{^{2}}}
\emergencystretch=2em
\begin{document}
\begin{center}\LARGE\bfseries Response to the Editor and Reviewers\end{center}\vspace{1em}

Submission ID: 20e66410-c472-42dc-ab99-957c24a43f96

Original title: Auditing observability bias in InSAR-derived subsidence exposure estimates with a reproducible geospatial workflow

Revised title: Quantifying how spatial allocation and sampling alter population-weighted InSAR support summaries

We thank the Editor and Reviewers for helping us sharpen the contribution and its evidential scope. The revised study makes a positive, testable claim: for a fixed population model, geographic domain and InSAR support field, spatial allocation and aggregation can materially change population-weighted support summaries. We quantify the size and direction of this effect, identify conditions that produce or reverse it, and test sensitivity to sampling and population models. We have removed the stronger interpretation of the original 3.77-million figure as a verified count of people exposed to subsidence in unobserved locations; this correction narrows the interpretation while preserving the supported methodological and empirical findings.

In addition to the requested reframing and point-by-point revisions, we conducted several further analyses to test where the conclusions hold and where they do not. These include a direct comparison of mean input-pair support with the final LiCSBAS quality mask and finite-domain bounds for unsupported pixels; controlled masking experiments that compare within-unit prevalence, centre sampling and a published unit-median-rate rule across missingness patterns and grid origins; crossed population-model and pair-sampling tests, plus scale/origin and fine-resolution allocation sensitivities; and Fresno benchmarks against Census blocks, CA-POP and ACS survey estimates. These additions test the robustness, transferability and limits of the revised claims. They do not provide physical validation in genuinely unsupported areas, and we state those limits alongside each result.

Location convention: manuscript and Supplementary Information page numbers refer to the clean PDFs. Manuscript changes are located by section, subsection and figure/table labels so they remain easy to find across line-wrap differences; the separately supplied numbered manuscript provides line numbers. Reviewer wording below is transcribed from the retained original comments. The marked copy shows changes since the local scientific baseline, not a tracked comparison with the submitted paper.

\section*{Editor}

\textbf{Comment.} In particular, the scientific gap and novelty of the study should be more clearly established and supported by a comprehensive and balanced review of the relevant literature. The rationale for the proposed workflow and its methodological choices should be strengthened and better positioned in relation to existing approaches. Further clarification and justification are also required regarding the interpretation and validation of the “hidden exposure”concept, the treatment of observability, the limitations of the GHSL data, and the selection of the InSAR datasets.

\textbf{Response.} We have restructured the Introduction around established InSAR quality control, published exposure practices and population areal interpolation, followed by three research questions. We also updated the literature comparison with recent work on nationwide population/building/infrastructure exposure (Xiong et al., 2026), building-level deformation-risk diagnosis (Yang et al., 2026), and uncertainty-aware coastal population-weighted relative sea-level rise (Oelsmann et al., 2026). These studies use broader inventories, target-level data or calibrated vertical-land-motion products; we distinguish their endpoints from our fixed-input allocation audit and do not imply that they validate our LOS observations. The revised gap does not depend on showing that earlier studies generally misused missing support. It asks how replacing a joint fine-scale population-support distribution with coarse summaries changes the reported quantity under fixed inputs, and under which conditions. The covariance decomposition explains the mechanism; the contribution is its quantitative evaluation across matched inputs, scales, origins and application settings, not a new mathematical identity.

The experimental response comprises matched-input scale/origin comparisons, controlled masking with known reference allocations, expanded sampling, alternative population products, a separate connected LiCSBAS experiment with excluded-edge residuals, and a Fresno application using actual published observation polygons. A further complete-interior analysis shows that the 1-km uniform-allocation difference remains 58.24-58.56\% after excluding boundary-crossing cells. This addresses a possible geometric alternative explanation. Internal rechecks and relocated-input checks are documented in Supplementary S11.

The supported finding is that the population-weighted support summary changes substantially when the same population total is redistributed uniformly within coarse cells rather than retained at native model resolution. This allocation-and-aggregation effect is quantified with matched-input tests, controlled references and a published observation-footprint case. It is not a count of verified residents over subsiding ground. GHSL/WorldPop disagreement is treated as population-model sensitivity, and the incomplete Iran branch is excluded from quantitative applications.

An added Supplementary S15 diagnostic evaluates the population--support covariance identity and its Cauchy--Schwarz bound across square sizes, grid origins and two population surfaces. The numerical identity closes within $1.22\times10^{-15}$ percentage points; we label this as an implementation check, not external validation or a confidence interval, and keep its pixel-centre sensitivity values separate from the direct-geometry primary estimates. The new crossed experiment tests the aggregation effect itself across population model and pair selection. With the common domain and population total fixed, uniform-minus-native differences are 55.38-55.64\% and 51.29-51.51\% for GHSL under 28 and 224 pairs, respectively; WorldPop gives 9.03-9.07\% and 8.64-8.68\%. All 16 model-sample-origin cases have a positive difference, while the magnitude is strongly model dependent. Four origins are sensitivity cases, not independent replicates. Pair number and composition change together. Locations: main manuscript p. 9, “Crossing population model and pair selection” paragraph; Supplementary S12 (Supplementary Information p. 13).

\textbf{Changes.} Updated frontier literature comparison, estimand definitions, boundary result, validation limits and historical source corrections in the main manuscript; Supplementary S7--S11, S15 and S18; additional matched-domain Fresno population-grid benchmark in Supplementary S19.

\subsection*{Code archive and submission files}

\textbf{Response.} The current public reviewer-reproduction package is the GitHub release \href{https://github.com/niubisile-wq/openinsar-observability-bias/releases/tag/v0.4.1}{v0.4.1}, fixed at commit \texttt{5bf98f19abbed981989f9c97ef9b118eea524c5b}. It contains the revised code, configurations, input manifests and checksums, processed analysis inputs, reference tables, environment records and reproduction instructions. A companion release asset supplies the six complete ACS source archives used for S24 and the two official Census source archives used for S20, with file-level hashes. Other provider-controlled source products remain available through the documented acquisition scripts and manifests; they are not claimed to be bundled. The earlier DOI \href{https://doi.org/10.5281/zenodo.21444768}{10.5281/zenodo.21444768} identifies v0.3.0 only and does not archive this revision. A new DOI-specific deposit has not yet been minted, so we identify the current revision by its public GitHub tag and commit; the DOI can be added to the Code Availability statement if and when it is assigned. The clean manuscript and the separately supplied local-baseline marked copy use the same current scientific text.

\section*{Reviewer 1}

\subsection*{R1.1}

\textbf{Comment (original wording).} The manuscript presents a valuable methodological contribution to the Earth science community by addressing an important but usually overlooked issue >>> how InSAR product support masks affect downstream exposure estimates.
There are some critical point that must be addressed.
1- The "Hidden Exposure" claim is not validated. The authors should justify such a claim in a scientific way as they address other issues.
The manuscript repeatedly emphasizes that it is not validating deformation in hidden cells, yet the central result (3.77 million people in "hidden exposure") is presented as if it carries decision-relevant weight.

\textbf{Response.} The reviewer identifies a real distinction between a spatial accounting estimate and a physically verified exposure count. We have corrected the interpretation accordingly. The 3,772,492.0508 value is reproducible from the submitted tables, but it sums whole coarse-cell populations selected by a center-majority rule; it cannot identify residents in unobserved subareas or establish subsidence there. We therefore no longer present it as a verified hidden-exposure count. This correction does not remove the paper’s central result: allocation and aggregation alter support summaries even when the domain, support field and population total are held fixed. We distinguish the conflicting archive branch and do not use the unconfirmed stored VLM scale as the primary population denominator.

Our primary estimand is mean area-pair support deficit. With geographic domain and support fixed, native GHSL allocation gives 3.020 million allocation units, whereas uniform redistribution within 1-km cells gives 4.763-4.770 million. These values quantify sensitivity to allocation relative to the selected native population model; they are not household counts. Controlled masks, analytical geometry, conservation checks, covariance decomposition and sign-reversing cases test the calculations against defined references. Synthetic deformation is explicitly synthetic, and its deterministic identification interval is not a confidence interval.

A further population cross-check compares the gridded models with official ACS 2017--2021 five-year survey estimates on 319 complete Fresno block groups. The aggregate 2020 Census count is within the exact ACS 90% replicate-based MOE, while GHSL and WorldPop raster totals differ from the ACS point estimate by +10.57% and +6.03%, respectively. We report this only as a coarser external survey benchmark: ACS is not an error-free truth set, source independence is unverified and block groups cannot validate household locations (Supplementary S24, pp. 26--27). This does not validate deformation in unsupported radar locations.

For the actual full-stack mask, we now also give a worst-case finite-domain bound: unsupported pixels contain 0.2944 percentage points of total allocation, so the LOS-class share interval has width 0.2944 pp at each of the -2, -5 and -10 mm yr$^{-1}$ thresholds. This makes explicit what remains unidentified without assigning deformation to unsupported locations (Supplementary S21, p. 23).

A separate LiCSBAS experiment processed 437 full inputs and 393 training inputs after 44 date-selected edges were excluded. After processing exclusions, 431 and 382 pairs respectively enter the inversions. All 44 excluded pairs are evaluable. The full quality mask accepts 97.02\% of pixels, including 83.34\% of the matched support-below-0.2 group. Thus low sampled support is not equivalent to final-product failure. The pooled excluded-edge RMSE is 0.984 mm, with 0.699/5.150 mm for training-mask pass/fail locations and 2.338/0.412 mm for the lowest/highest training-support bins. Time-series differences, training linear velocity and zero relative displacement are compared on identical evaluable observations. The strata use training data only.

The LiCSBAS experiment adds a distinct processing-chain result: low sampled support is not equivalent to failure of the final product’s quality mask, and excluded-edge residuals provide an internal reconstruction check. These are held-out edges between dates already represented in training, so shared dates, reference handling and atmospheric signals limit independence. The residuals do not establish out-of-time predictive accuracy, GNSS/leveling accuracy or motion in permanently missing areas. Published Bangkok validation studies and the SBKK support check provide context, not validation of the Chao Phraya allocation estimate.



We now directly pair the 437-input-pair mean-support field and the final full-stack quality mask with the same population field on the actual LiCSBAS grid. Independent source-cell intersections confirm that 1-km uniform allocation increases the deficit by 3.066--3.189 percentage points for mean input support and by 0.108--0.117 points for the final mask across four grid origins. Their native deficits (3.4539\% and 0.2944\%) refer to different support definitions, so we report them separately rather than treating one as validation of the other. This confirms that an allocation effect remains for the actual final quality mask, while showing that it is much smaller than for mean input-pair support. The added experiment is documented in Supplementary S13; it still does not identify motion in invalid pixels.

We also decomposed these two effects spatially on fixed 1-km and 5-km UTM grids. At the primary 1-km scale, the uniform-minus-native difference is 3.194 percentage points for mean input support and 0.103 points for the final mask, consistent in scale with the exact-intersection estimates above. For the final mask, area-uniform and native-population support labels differ at the descriptive 50\% cutoff in 17 1-km units containing 2,229 GHSL allocation units; the corresponding input-support comparison has 150 units and 139,903 allocation units. We therefore report the final-mask allocation effect as small and make the support definition explicit; the cutoff is not a hazard threshold, and the comparison does not establish motion in invalid pixels (Supplementary S18).



A second added test uses artificial masks on actual fitted LOS velocities within the original final-valid domain. At a predeclared $-5$ mm yr$^{-1}$ LOS cutoff, 30\% removal and nominal 1.11-km blocks, the local-prevalence estimator's mean relative error was +1.1\% for random masks, +5.2\% for clustered masks, $-30.1\%$ when high-spread pixels were preferentially removed and +79.1\% when low-spread pixels were removed. Nearest-centre sampling was +99.7\% biased under random masking. This result demonstrates mechanism dependence under a known artificial reference; it does not impute the genuinely invalid pixels or establish a hazard threshold. The full factorial appears in Supplementary S14.

We additionally intersected actual 2020 Census block geometries and population counts with the Fresno DWR observation polygons. On 6,788 complete blocks, the population-weighted area-uniform support share is 99.19\%; classifying blocks by whether at least half of block area is inside the observation footprint assigns 99.41\% of Census population to the majority-observed class (Supplementary S16). We then used the separate DWR/TRE ALTAMIRA annual vertical-rate rasters to implement the published median-rate Census-block population-transfer step. In the 2015--2021 temporal median of six annual products, valid rates cover 6,129 blocks containing 501,411 people (95.70\% of complete-block Census population); 65.23\% of that covered population falls in blocks with median rates below $-5$ mm yr$^{-1}$. The analysis now provides actual block-median rate classes, while remaining distinct from our original LOS field: DWR rates are interpolated vertical products, and their processing lineage overlaps the DWR observation-footprint product. The endpoint-rate sensitivity also yields materially different class shares, so these results are a transferability demonstration, not independent validation or a like-for-like replication of our InSAR stack. Census counts are differentially private, and GHSL source independence is unverified. The full methods, temporal sensitivity and per-block results are in Supplementary S17. Within the 2015--2021 epoch, a six-window persistence analysis further shows that 95.16\% of residents are in blocks that cross the $-5$ mm yr$^{-1}$ threshold in at least one of six annual windows (95.06\% in one to five windows plus 0.10\% in all six) (Supplementary Table S17). We also added an exploratory transition matrix on the 6,129 blocks observed in both periods: holding 2020 Census counts fixed, 60.88\% of residents change median-rate class between 2015--2021 and the 2025--2026 single-year layer. This comparison is explicitly labeled as temporal-product sensitivity, not contemporary exposure or long-term hazard change (Supplementary Table S16).

To quantify the resolution limits of the population surfaces used in the transfer analyses, we added a matched-domain comparison of GHSL and WorldPop 1-km grids against these same 2020 Census block totals. Both normalized grids outperform a baseline that allocates population only by block land area, but block-level WAPE remains 59.37\% for GHSL and 57.20\% for WorldPop. We therefore interpret this as a coarse model-agreement check rather than household-scale validation: Census counts are differentially private, GHSL has census-derived inputs, and exact source overlap is unresolved. It does not validate exposed-person counts. The experiment and per-block data are reported in Supplementary S19 and retained with checksums in the analysis bundle.

\textbf{Changes and evidence.} main manuscript pp. 9-10, “Transfer to actual published support” Results and the Fresno block-transfer and population-grid benchmark paragraphs; Supplementary S16-S19 and S24 (Supplementary Information pp. 17-21 and 26); main manuscript and Supplementary S1-S4, S7, S11 for the revised estimand, validation limits and controlled-reference checks.

\subsection*{R1.2}

\textbf{Comment (original wording).} P 10: The classification of cells as "majority observable" vs. "not-majority-observable" simplifies a continuous problem. I think binary threshold is unsatisfying the goals and scopes you defined before.

\textbf{Response.} We have made support continuous in the primary analysis and report population-, built-surface- and area-weighted distributions over [0,1], rather than making the majority split the principal result. The mean deficit equals the integral of the weighted empirical support distribution; we verify this from the exact step distribution. Cutoffs 0.2, 0.3 and 0.4 and unwrapped validity alone remain sensitivity definitions.

The mean does not distinguish spatially persistent missingness from variable pass frequency, so these are not relabeled as the same object. Nor is the 28-pair diagnostic treated as a complete inversion network: it has 54 dates and 26 components. We report expanded 28/56/112/224-pair samples, temporal strata and 100 conditional draws per sample size, while keeping the separate connected processing experiment distinct. The original and nested 28-pair samples differ and are now explicitly identified.

\textbf{Changes and evidence.} main manuscript pp. 3-5, Materials and methods and the opening Results; Supplementary S2 and S5, Figs. S1-S2 and the acquisition-network table. Main Fig. 2 shows the full support distributions.

\subsection*{R1.3}

\textbf{Comment (original wording).} P9-10: "GHSL layers are used because they provide consistent public support for population and built-up surface across domains." >>> GHSL data limitations are understated and the authors should improve this statements.

\textbf{Response.} Population-model sensitivity is now quantified directly. On the common valid domain the q=0.3 deficit is 13.37\% for GHSL 2020 and 23.20\% for WorldPop 2020. Equalizing the population totals still gives 2.989 versus 5.186 million units, so total-population differences alone do not explain the contrast. We report each model’s domain, joint coverage, excluded population and normalized results.

GHSL 2015/2020 epochs are compared separately at fixed 30-arc-second resolution; native-model allocation is compared across reporting scales. Fractional WorldCover summaries preserve an explicit uniformity assumption within GHSL pixels. GHSL and WorldPop remain census/model-based allocations, not local demographic truth. GHSL population and built-surface layers share information, so a built-weighted comparator is not independent validation. Built surface is not occupied floor area; population is not vulnerability, loss or damage. Model differences are not calibrated error bars.

The new crossed experiment tests the aggregation effect itself across population model and pair selection. With the common domain and population total fixed, uniform-minus-native differences are 55.38-55.64\% and 51.29-51.51\% for GHSL under 28 and 224 pairs, respectively; WorldPop gives 9.03-9.07\% and 8.64-8.68\%. All 16 model-sample-origin cases have a positive difference, while the magnitude is strongly model dependent. Four origins are sensitivity cases, not independent replicates. Pair number and composition change together. Locations: main manuscript p. 9, “Crossing population model and pair selection” paragraph; Supplementary S12 (Supplementary Information p. 13).

\textbf{Changes and evidence.} main manuscript pp. 3-5 and 9-12, Methods, matched-processing and Discussion sections; main Fig. 5 and Supplementary S6, including all common-domain denominators.

\subsection*{R1.4}

\textbf{Comment (original wording).} there are also some minor issues that must be double-checked by the authors.
I recommend major revisions with particular attention to clarifying the weak points and the policy-relevant interpretation of the results too.

\textbf{Response.} We have checked the scientific quantities, units, dates, sample counts, figure captions, source attribution and cross-references. The revisions distinguish the original 3.772492-million center-majority total, an early revision estimand and the current 3.020-million mean deficit. They also distinguish source-product coverage, sampled pair support and final LiCSBAS quality. The corrected Iran attribution, unconfirmed VLM scale, population-model dependence and temporal mismatches are explicitly retained.

We remove policy claims about a verified population hidden from subsidence assessment and do not propose a validated measurement-priority ranking. The practical recommendation is limited to transparent reporting of the domain, support definition, allocation assumptions, sensitivities and unresolved information. Internal numerical checks and PDF review support consistency; they do not add external physical evidence.

\textbf{Changes and evidence.} main manuscript p. 12, Discussion; Supplementary S1 and S11 (Supplementary Information p. 11).

\section*{Reviewer 2}

The manuscript "Auditing observability bias in InSAR-derived subsidence exposure estimates with a reproducible geospatial workflow" by Zixuan Liu and Wei Xiong discusses a very important topic regarding InSAR products and how they might translate into misleading exposure estimates. While the topic is relevant and the issues raised by the manuscript are valid, I have major concerns about the structure of the paper.

\subsection*{R2.1}

\textbf{Comment (original wording).} 1- The problem formulation gives too little credit to established InSAR practices. It discusses the problem in the Introduction. For example:

"exposure accounting can convert this missingness into a denominator change."

"A processing team may publish a careful product with valid-support masks, while a downstream analyst may treat the remaining raster or point layer as a complete hazard domain."

"but they rarely audit whether non-observed radar support is being treated as absent"

I find the concern relevant and valid. But the literature review does not systematically show that previous InSAR exposure studies have failed to address areas with missing information. The manuscript claims that "When the resulting deformation map is intersected with population, built-up land or infrastructure layers, analysts often remove cells that are not valid in the radar product before exposure is summarized," without providing evidence from studies that follow this practice. In my opinion, this is the main weakness of the manuscript. If the issue is misuse of correctly documented InSAR products by non-specialist users, this should be stated explicitly and supported by multiple examples of studies where such problems occur, together with references to state-of-the-art studies that have tried to address the issue. In its current form, the manuscript sometimes reads as if the InSAR community has failed to recognize this limitation. However, it is well accepted in the community that InSAR observations are opportunistic by nature. In some cases, InSAR products are accompanied by quantitative information on the validity and quality of the observations. For example, LiCSBAS provides average coherence, the percentage of valid unwrapped pixels, etc. along with velocity and time-series products. The manuscript acknowledges that experienced InSAR users understand and correctly interpret areas where InSAR does not provide information. Therefore, the scientific gap should be demonstrated specifically at the point where InSAR products are used in downstream applications. Without enough evidence from previous studies, it is difficult to determine whether the paper addresses a genuine unresolved research problem.

\textbf{Response.} We have revised the framing to recognize established InSAR practice, including PS, SBAS, distributed-scatterer, LiCSBAS and MintPy methods and their quality diagnostics. The literature review describes what relevant exposure papers actually report and does not infer widespread misuse from selected examples. Our research question remains substantive without that prevalence claim: downstream summaries can change when fine-scale population-support structure is aggregated, even when the input support and population totals are unchanged.

Examples include Blackwell et al.'s stated uniform-within-community allocation and GNSS comparisons; Fernandez-Torres et al.'s coherence screening and GPS checks; Ao et al.'s urban population shares above rate thresholds; Ohenhen et al.'s census-block median-rate assignment and comparison with 154 GNSS stations; and Murray et al.'s explicit limitations for motion between buildings. We describe these as relevant practices, not demonstrations of misuse. Published Bangkok ground comparisons are summarized in the Introduction, with detailed station/time/support limitations moved to Supplementary S8.

The empirical question is: for a fixed population model, domain and support field, how much does spatial aggregation change the population-weighted support summary, under which conditions, and in which direction? Matched-input comparisons, complete-interior results, controlled sign reversals and model/sampling sensitivity answer it directly. This is a contribution about quantifying downstream aggregation effects, not an allegation that earlier authors erred. We distinguish LiCSAR sampled support from the WabInSAR source product’s validity mask; Fresno tests transfer to a published observation footprint with a different support definition.

The crossed experiment further demonstrates that this contribution is not confined to the primary GHSL calculation: its direction persists for WorldPop and expanded pair sampling, with a much smaller magnitude for WorldPop. We now explicitly restrict the approximately 58\% value to the original primary case (main manuscript p. 9, “Crossing population model and pair selection” paragraph; Supplementary S12, Supplementary Information p. 13).

We have added a direct controlled benchmark of a published census-unit median-rate assignment rule against within-unit observed-population prevalence and centre sampling. Across 100 mask replicates, four fixed grid origins and four missingness mechanisms, the median-rate rule has about +1.9 pp mean signed error under random and clustered removal, while within-unit prevalence is +0.08 and +0.14 pp, respectively; under low-spread removal the median-rate rule is closer, and under high-spread removal all three underestimate. These are artificial-mask design sensitivities on originally valid LOS data, not evidence about genuine missing pixels (Supplementary S22, p. 24). This comparison both strengthens the downstream-method rationale and limits any claim of universal superiority.

\textbf{Changes and evidence.} main manuscript pp. 1-2, Introduction, and p. 12, Discussion; Supplementary S8 (Supplementary Information p. 8) and S10 (Supplementary Information p. 10).

\subsection*{R2.2}

\textbf{Comment (original wording).} 2- The manuscript explains how the proposed workflow works, but it does not adequately explain why this particular method is chosen. State-of-the-art methods and alternative approaches should be discussed and compared with the proposed method. The Materials and Methods section also contains very few references to previous methodological works. As a result, it is difficult to determine how the proposed approach relates to existing methods and why the specific choices made in the workflow are preferable or necessary.

\textbf{Response.} The revised Methods and Supplementary S10 position the workflow against established alternatives, with methodological references including Eicher and Brewer’s dasymetric comparison and Mennis’s areal/ancillary allocation work. We compare inputs, assumptions, interpretation and evaluation for center sampling, uniform areal allocation, native-model allocation, built weighting, observed-domain overlays, processing diagnostics and deformation imputation/additional sensors.

For methods estimating the same quantity, support, domain and population totals are fixed. Four scales and four origins test center and uniform allocation against the native model, while local errors and a covariance decomposition explain the discrepancy. The new complete-interior result remains 58.24-58.56\%, so the primary effect is not explained by boundary-crossing cells. The controlled-mask experiments retain both signs of error, local cancellation and small aggregate-bias cases. Fresno repeats the comparison using actual polygons; the small outside-support denominator is accompanied by an error normalized to the whole AOI population.

Native allocation is useful when retaining the selected model's fine structure is appropriate; it does not validate that structure. Uniform allocation is adequate when the relevant covariance is negligible at the required scale. Built weights add assumptions and share information with GHSL. Imputation, sensor fusion and ground measurements address missing deformation under additional observations or assumptions; we do not claim numerical superiority over such different-target methods. The rationale for our protocol is conserving and exposing the consequences of known inputs, with quantified applicability conditions.

To test method sensitivity across inputs, the additional 2-by-2 model/pair comparison repeats the native-versus-uniform calculation under identical domains and totals. It preserves the population-support product before projection and reports signed differences against both the native deficit and the whole population. The 25-m integration diagnostic changes the result by at most 0.000436 percentage points of total population. The full 16-row table and numerical checks are supplied in Supplementary S12 (Supplementary Information p. 13).



We add one external population benchmark at the Fresno block-group scale. Against ACS 2017--2021 estimates, the 2020 Census P1 total differs by $-0.36\%$ (inside the aggregate ACS 90\% MOE), while the GHSL and WorldPop 1-km totals differ by $+10.57\%$ and $+6.03\%$. The exact aggregate ACS uncertainty uses the 80 published variance replicates. This identifies model disagreement at the reporting-unit scale and does not establish which raster is correct at household scale (Supplementary S24, pp. 26--27).

We also report worst-case finite-domain bounds for LOS threshold-class population allocation under the actual full-stack mask. The unsupported allocation mass is 0.2944 percentage points of the fixed domain, giving class-share widths of 0.2944 pp at the -2, -5 and -10 mm yr$^{-1}$ thresholds (Supplementary S21, p. 23). These bounds make the unresolved allocation explicit; they do not infer motion in unsupported pixels. A prespecified-tolerance interpretation of the scale/origin covariance bound is added in Supplementary S23, p. 25.

We added a paired LiCSBAS analysis that holds the population field, processing grid and coarse allocation constant while separately evaluating mean support across 437 inputs and the final full-stack quality mask. Direct intersections of source cells with the 1-km reporting grid verify the geometry independently of raster reprojection. The measured effects differ materially by support definition, demonstrating why alternative methods must be compared on the same estimand and why sampled input support should not be used as a stand-in for the final product mask. We report the four-origin sensitivity and full numerical audit in Supplementary S13.



We added a fine-resolution population-allocation sensitivity using official 2020 Census P1 block totals and the CA-POP 100-m dasymetric surface in Fresno. Normalizing both GHSL and CA-POP within the same blocks fixes their totals and isolates within-block allocation assumptions against the same DWR polygon-union support. The outside-support deficit is 2,016.3 Census allocation units for GHSL and 1,268.4 for CA-POP (a difference of 0.143 percentage points of the shared total); the direction varies across tracts. Because CA-POP is also built from these Census block counts, this is model sensitivity, not independent population validation. We quantified spatial dependence among the block-level unsupported fractions using binary queen-contiguity Moran's $I$: the CA-POP-minus-GHSL difference field has $I=0.399$ on 17,083 neighboring block pairs. This is reported descriptively for the fixed domain without a sampling inference. We also tested 250--2,000-m grids at four origins for center and uniform allocation, reporting local absolute errors, ranking overlap and a Cauchy--Schwarz bound; ranking overlap declines at coarser scales and center sampling varies with origin. The new tests support explicit scale/origin reporting rather than a universal correction factor (Supplementary S20, p. 22; main manuscript Fresno application).

The observed-LOS masking experiment also evaluates the proposed allocation choices against a known class-population reference created by masking pixels with measured LOS rates. Its results show good aggregate behavior for within-block prevalence under random masks but larger, sign-varying errors under spatially clustered or quality-selective masks; nearest-centre sampling can be substantially biased even under random masks. This strengthens the comparison with an observed processing product while keeping the estimand explicit: only hidden values inside the originally valid domain are tested. The experiment does not supply evidence about permanent no-data areas.

\textbf{Changes and evidence.} main manuscript pp. 3-5, Methods and Results; p. 7, Controlled missingness; pp. 9-10, scale, population-model and Fresno benchmark results; Supplementary S3-S4, S9-S11 and S24 (Supplementary Information pp. 10-11 and 26); main Figs. 3-4 and 7.

\subsection*{R2.3}

\textbf{Comment (original wording).} 3- The "Iran public InSAR products" is based on relatively old study from 2008. More recent nationwide studies of subsidence in Iran are available, including Haghighi and Motagh, 2024 and Payne et al., 2025. InSAR datasets form both of these studies are publicly available. The manuscript should explain why an older dataset was selected instead of these newer products. I think the use of the 2008 dataset does not invalidate the workflow, but the it should be justified.

\textbf{Response.} The source audit changes the premise of the original selection question: checksum comparison identifies the archived rate and mask as the 2014-2020 product accompanying Haghshenas Haghighi and Motagh (2024), not a 2008 nationwide rate product. The archived seasonal-amplitude raster does not match the current companion checksum and is excluded. Because the earlier Iran branch tested input parsing but did not complete a matched-input application, we do not present Iran as a quantitative case or claim a new/old product comparison.

The revised source table distinguishes Haghshenas Haghighi and Motagh's national annual-rate raster and rate-derived mask from Payne et al.'s 2014-2022 selected-region vertical/east-west velocities and uncertainties. Those summary layers are not pair-specific valid-unwrapped/coherence masks. Underlying LiCSAR interferograms are separately available; we do not claim that a future reprocessing is impossible. We have not performed a new/old Iran product comparison. The quantitative applications are now Chao Phraya and the explicitly dated Fresno DWR observation polygons.

\textbf{Changes and evidence.} main manuscript p. 12, Discussion; pp. 5 and 9, matched-input and scale-sensitivity Results; Supplementary S1 and S8 (Supplementary Information p. 8).
\end{document}
